% Źródła: materialy/2026-10-27/zmienna-losowa.png i IMG_7006--IMG_7013.jpeg.
\chapter{Zmienne losowe}
\label{chap:zmienne-losowe}

\section{Definicja i mierzalność}
\label{sec:definicja-zmiennej-losowej}

Niech $(\Omega,\F,\PP)$ będzie przestrzenią probabilistyczną.

\begin{definicja}
\label{def:zmienna-losowa}
\emph{Zmienną losową o wartościach w $\R$} na przestrzeni
$(\Omega,\F,\PP)$ nazywamy funkcję $X\colon\Omega\to\R$ taką, że
\[
  X^{-1}(A)=\{\omega\in\Omega:X(\omega)\in A\}\in\F
  \qquad\text{dla każdego }A\in\B(\R).
\]
Tutaj $\B(\R)$ oznacza $\sigma$-ciało zbiorów borelowskich prostej.
\end{definicja}

\begin{uwaga}
Analogicznie definiujemy zmienną losową $X\colon\Omega\to\R^d$,
wymagając, aby $X^{-1}(A)\in\F$ dla każdego $A\in\B(\R^d)$.
Taką zmienną nazywamy \emph{wektorem losowym}.
\end{uwaga}

\begin{uwaga}
Zmienna losowa jest funkcją mierzalną.
\end{uwaga}

\begin{uwaga}
Aby funkcja $X\colon\Omega\to\R$ była zmienną losową, wystarczy,
żeby
\[
  \{\omega\in\Omega:X(\omega)>a\}\in\F
  \qquad\text{dla każdego }a\in\R.
\]
\end{uwaga}

\begin{uwaga}[Zapis prawdopodobieństw]
Dla zbioru borelowskiego $A\subseteq\R$ stosujemy skrót
\[
  \PP(X\in A)=\PP\bigl(\{\omega\in\Omega:X(\omega)\in A\}\bigr).
\]
\end{uwaga}

\begin{uwaga}[Mierzalność funkcji]
Funkcja $f\colon(E,\mathcal E)\to(S,\mathcal S)$ jest
\emph{mierzalna}, jeżeli
\[
  f^{-1}(A)\in\mathcal E
  \qquad\text{dla każdego }A\in\mathcal S.
\]
\end{uwaga}

\begin{fakt}
\label{fakt:zlozenie-zmiennej-losowej}
Jeżeli $X\colon\Omega\to\R^d$ jest wektorem losowym, a funkcja
$\varphi\colon\R^d\to\R^m$ jest borelowska, to
\[
  \varphi(X)=\varphi\circ X\colon\Omega\to\R^m
\]
jest wektorem losowym (dla $m=1$ --- zmienną losową rzeczywistą).
\end{fakt}

\begin{uwaga}
Funkcja $\varphi\colon\R^d\to\R^m$ jest \emph{borelowska}, jeżeli
\[
  \varphi^{-1}(A)\in\B(\R^d)
  \qquad\text{dla każdego }A\in\B(\R^m).
\]
\end{uwaga}

\section{\texorpdfstring{$\sigma$}{sigma}-ciało generowane przez zmienną losową}
\label{sec:sigma-zmiennej-losowej}

\begin{definicja}
\label{def:sigma-zmiennej-losowej}
Niech $X\colon\Omega\to\R$ będzie zmienną losową.
Najmniejsze $\sigma$-ciało $\mathcal G\subseteq\F$, względem którego
$X$ jest mierzalna, nazywamy \emph{$\sigma$-ciałem generowanym przez $X$}
i oznaczamy przez $\sigma(X)$.
\end{definicja}

\begin{fakt}
\label{fakt:postac-sigma-x}
Dla zmiennej losowej $X\colon\Omega\to\R$ mamy
\[
  \sigma(X)=\{X^{-1}(B):B\in\B(\R)\}.
\]
\end{fakt}

\begin{przyklad}[Zmienna przyjmująca dwie wartości]
\label{przykl:sigma-dwie-wartosci}
Niech $A\in\F$ oraz $a,b\in\R$, $a\ne b$. Rozważmy zmienną losową
\[
  X(\omega)=a\mathbf{1}_A(\omega)+b\mathbf{1}_{A^c}(\omega),
  \qquad\omega\in\Omega,
\]
gdzie $\mathbf{1}_A$ jest funkcją wskaźnikową zbioru $A$.
Dla $B\in\B(\R)$ otrzymujemy
\[
  X^{-1}(B)=
  \begin{cases}
    \varnothing,&a\notin B,\ b\notin B,\\
    A,&a\in B,\ b\notin B,\\
    A^c,&a\notin B,\ b\in B,\\
    \Omega,&a\in B,\ b\in B.
  \end{cases}
\]
Zatem
\[
  \sigma(X)=\{\varnothing,A,A^c,\Omega\}.
\]
\end{przyklad}

\section{Niezależność zmiennych losowych}
\label{sec:niezaleznosc-zmiennych-losowych}

\begin{definicja}
\label{def:niezaleznosc-zmiennych-losowych}
Zmienne losowe $X,Y\colon\Omega\to\R$ są \emph{niezależne}, jeżeli
dla każdych $A,B\in\B(\R)$ zachodzi
\[
  \PP(X\in A,\,Y\in B)=\PP(X\in A)\PP(Y\in B).
\]
Równoważnie: $\sigma$-ciała $\sigma(X)$ i $\sigma(Y)$ są niezależne.
\end{definicja}

\begin{definicja}
\label{def:niezaleznosc-rodziny-zmiennych}
Zmienne losowe $X_1,\ldots,X_n$ są \emph{niezależne}, jeżeli
$\sigma$-ciała $\sigma(X_1)$, $\ldots$, $\sigma(X_n)$ są niezależne.
Ogólniej, dla dowolnego zbioru indeksów $T$ rodzina zmiennych losowych
$(X_t)_{t\in T}$ jest \emph{niezależna}, jeżeli rodzina
$\bigl(\sigma(X_t)\bigr)_{t\in T}$ jest niezależna.
\end{definicja}

\begin{przyklad}[Niezależne funkcje wskaźnikowe]
\label{przykl:niezalezne-wskaznikowe}
Rozważmy przestrzeń probabilistyczną
\[
  (\Omega,\F,\PP)=\bigl([0,1],\B([0,1]),\lambda\bigr),
\]
gdzie $\lambda$ oznacza miarę Lebesgue'a na odcinku $[0,1]$.
Niech
\[
  A=\left[0,\frac12\right),
  \qquad
  B=\left[0,\frac14\right)\cup\left[\frac12,\frac34\right),
  \qquad X=\mathbf{1}_A,\quad Y=\mathbf{1}_B.
\]
Wówczas
\[
  \sigma(X)=\{\varnothing,A,A^c,\Omega\},
  \qquad
  \sigma(Y)=\{\varnothing,B,B^c,\Omega\}.
\]
Ponadto
\[
  \PP(A\cap B)
    =\lambda\!\left(\left[0,\frac14\right)\right)
    =\frac14=\frac12\cdot\frac12=\PP(A)\PP(B).
\]
Zdarzenia $A,B$ są niezależne, a z własności niezależności
(fakt~\ref{fakt:wlasnosci-niezaleznosci}) to samo dotyczy ich dopełnień.
Zatem $\sigma$-ciała $\sigma(X)$ i $\sigma(Y)$ są niezależne,
czyli zmienne losowe $X,Y$ są niezależne.
\end{przyklad}

\begin{przyklad}[Funkcje Rademachera]
\label{przykl:rademacher}
Na tej samej przestrzeni probabilistycznej funkcje
\[
  X_n(x)=\operatorname{sgn}\!\bigl(\sin(2^n\pi x)\bigr),
  \qquad x\in[0,1],\quad n=1,2,\ldots,
\]
tworzą rodzinę niezależnych zmiennych losowych.
Zera sinusów występują w punktach dyadycznych, w tym na końcach odcinka;
zbiór tych punktów ma miarę zero, więc wartości funkcji w tych punktach
nie wpływają na niezależność.
\end{przyklad}

\begin{fakt}
\label{fakt:zmienna-niezalezna-od-siebie}
Zmienna losowa $X\colon\Omega\to\R$ jest niezależna od samej siebie
wtedy i tylko wtedy, gdy istnieje $c\in\R$ takie, że
\[
  \PP(X=c)=1.
\]
\end{fakt}

\section{Układy \texorpdfstring{$\pi$ i $\lambda$}{pi i lambda}}
\label{sec:pi-lambda}

\begin{definicja}
\label{def:pi-uklad}
Rodzinę zbiorów $\mathcal P\subseteq2^\Omega$ nazywamy
\emph{$\pi$-układem}, jeżeli
\[
  A\cap B\in\mathcal P
  \qquad\text{dla wszystkich }A,B\in\mathcal P.
\]
\end{definicja}

\begin{przyklad}
Rodzina
\[
  \mathcal P=\{[a,b):a,b\in\R\}
\]
jest $\pi$-układem na $\R$. Przyjmujemy przy tym $[a,b)=\varnothing$
dla $a\geq b$.
\end{przyklad}

\begin{definicja}
\label{def:lambda-uklad}
Rodzinę zbiorów $\mathcal L\subseteq2^\Omega$ nazywamy
\emph{$\lambda$-układem}, jeżeli spełnia następujące warunki:
\begin{enumerate}
  \item $\Omega\in\mathcal L$.
  \item Jeżeli $A,B\in\mathcal L$ i $A\subseteq B$, to
    $B\setminus A\in\mathcal L$.
  \item Jeżeli $A_1,A_2,\ldots\in\mathcal L$ oraz
    $A_n\subseteq A_{n+1}$ dla każdego $n\geq1$, to
    \[
      \bigcup_{n=1}^{\infty}A_n\in\mathcal L.
    \]
\end{enumerate}
\end{definicja}

\begin{fakt}
\label{fakt:pi-lambda-sigma}
Jeżeli rodzina $\F$ jest jednocześnie $\pi$-układem i $\lambda$-układem,
to jest $\sigma$-ciałem.
\end{fakt}

\begin{twierdzenie}[O układach $\pi$ i $\lambda$]
\label{tw:pi-lambda}
Jeżeli $\lambda$-układ $\mathcal L$ zawiera $\pi$-układ $\mathcal P$,
to
\[
  \sigma(\mathcal P)\subseteq\mathcal L.
\]
\end{twierdzenie}

\subsection{Sposób stosowania twierdzenia}
\label{subsec:stosowanie-pi-lambda}

Chcąc udowodnić pewną własność $(\ast)$ dla każdego $A\in\F$,
możemy postępować następująco:
\begin{enumerate}
  \item Znajdujemy $\pi$-układ $\mathcal P$ taki, że
    $\sigma(\mathcal P)=\F$, i sprawdzamy własność $(\ast)$
    dla każdego zbioru z $\mathcal P$.
  \item Definiujemy rodzinę
    \[
      \mathcal L=\{A\in\F:\text{własność }(\ast)\text{ zachodzi dla }A\}.
    \]
  \item Wykazujemy, że $\mathcal L$ jest $\lambda$-układem.
  \item Ponieważ $\mathcal P\subseteq\mathcal L$, z
    twierdzenia~\ref{tw:pi-lambda} otrzymujemy
    \[
      \F=\sigma(\mathcal P)\subseteq\mathcal L.
    \]
    Własność $(\ast)$ zachodzi więc dla każdego $A\in\F$.
\end{enumerate}
